% GENERATED FILE. Edit canonical lessons, then run npm run generate:books.
% canonical-source-hash: fnv1a64:af74ebe721f90991
% canonical-lessons: UR-C475-daimi, UR-C475-dama, UR-C475-falij, UR-C475-kainsar, UR-C475-mubtala

\chapter{Long Illnesses}
\label{ch:ur-long-illnesses}
\hlchaptermodality{475}

\begin{chapteropening}
\textbf{By the end of this chapter:} \emph{I can talk about an illness that lasts.}

Complete the last lesson of chapter 475: I can talk about an illness that lasts.
\end{chapteropening}

\section[\texorpdfstring{\ur{دائمی} (dāimī)}{دائمی (dāimī)}]{\ur{دائمی} (dāimī) — chronic; lasting}

\label{lesson:UR-C475-daimi}



\begin{quote}
Before the new one: say the Urdu for measles, then the Urdu for smallpox.
\end{quote}

\subsection*{You'll want to know: \ur{دائمی}}
\textbf{\ur{دائمی}} — \emph{dāimī} — \textquotedblleft{}chronic; lasting\textquotedblright{}.

\begin{quote}
\ur{دادی} \ur{کو} \ur{دائمی} \ur{کھانسی} \ur{ہے۔}
\end{quote}

\begin{quote}
\emph{dādī ko dāimī khāṅsī hai.}
\end{quote}

\begin{quote}
My grandmother has a chronic cough.
\end{quote}

\begin{cousinweb}
\textbf{\ur{دائمی}} (\emph{dāimī}) is usually traced, through Persian, to Arabic \emph{dāʾimī}, \textquotedblleft{}lasting, perpetual\textquotedblright{}.
\end{cousinweb}

\begin{grammarlens}[title={What you've built}]
One more word for describing what you lived and felt.
\end{grammarlens}

\subsection*{Your turn}
\begin{itemize}
\raggedright
  \item \emph{Say it:} \emph{dāimī}
  \item \emph{Say it:} \emph{dāimī}, once more
  \item \emph{Say it:} \emph{dādī ko dāimī khāṅsī hai.} — My grandmother has a chronic cough.
  \item \emph{Recall:} say the Urdu for measles, then the Urdu for smallpox, then say \emph{dāimī} again
\end{itemize}

\begin{tcolorbox}[breakable,colback=teal!4,colframe=teal!35!black,title={Before you move on}]
What does \ur{دائمی} mean? (\textquotedblleft{}Chronic; lasting\textquotedblright{}.) Say it once more.
\end{tcolorbox}
\section[\texorpdfstring{\ur{دمہ} (damā)}{دمہ (damā)}]{\ur{دمہ} (damā) — asthma}

\label{lesson:UR-C475-dama}



\begin{quote}
Before the new one: say the Urdu for smallpox, then the Urdu for chronic.
\end{quote}

\subsection*{You'll want to know: \ur{دمہ}}
\textbf{\ur{دمہ}} — \emph{damā} — \textquotedblleft{}asthma\textquotedblright{}.

\begin{quote}
\ur{دمے} \ur{کے} \ur{مریض} \ur{کو} \ur{سانس} \ur{لینے} \ur{میں} \ur{مشکل} \ur{ہوتی} \ur{ہے۔}
\end{quote}

\begin{quote}
\emph{dame ke marīz ko sāṅs lene meṅ mushkil hotī hai.}
\end{quote}

\begin{quote}
Someone with asthma finds it hard to breathe.
\end{quote}

Masculine.

\begin{cousinweb}
\textbf{\ur{دمہ}} (\emph{damā}) is usually traced to Persian \emph{dama}, \textquotedblleft{}asthma\textquotedblright{}, beside \emph{dam}, \textquotedblleft{}breath\textquotedblright{} (\textbf{\ur{دم}}).
\end{cousinweb}

\begin{grammarlens}[title={What you've built}]
One more word for describing what you lived and felt.
\end{grammarlens}

\subsection*{Your turn}
\begin{itemize}
\raggedright
  \item \emph{Say it:} \emph{damā}
  \item \emph{Say it:} \emph{damā}, once more
  \item \emph{Say it:} \emph{dame ke marīz ko sāṅs lene meṅ mushkil hotī hai.} — Someone with asthma finds it hard to breathe.
  \item \emph{Recall:} say the Urdu for smallpox, then the Urdu for chronic, then say \emph{damā} again
\end{itemize}

\begin{tcolorbox}[breakable,colback=teal!4,colframe=teal!35!black,title={Before you move on}]
What does \ur{دمہ} mean? (\textquotedblleft{}Asthma\textquotedblright{}.) Say it once more.
\end{tcolorbox}
\section[\texorpdfstring{\ur{فالج} (fālij)}{فالج (fālij)}]{\ur{فالج} (fālij) — paralysis; a stroke}

\label{lesson:UR-C475-falij}



\begin{quote}
Before the new one: say the Urdu for chronic, then the Urdu for asthma.
\end{quote}

\subsection*{You'll want to know: \ur{فالج}}
\textbf{\ur{فالج}} — \emph{fālij} — \textquotedblleft{}paralysis; a stroke\textquotedblright{}.

\begin{quote}
\ur{دادا} \ur{کو} \ur{فالج} \ur{ہوا} \ur{تھا۔}
\end{quote}

\begin{quote}
\emph{dādā ko fālij huā thā.}
\end{quote}

\begin{quote}
My grandfather had had a stroke.
\end{quote}

Masculine.

\begin{cousinweb}
\textbf{\ur{فالج}} (\emph{fālij}) is usually traced, through Persian, to Arabic \emph{fālij}, \textquotedblleft{}paralysis\textquotedblright{}.
\end{cousinweb}

\begin{grammarlens}[title={What you've built}]
One more word for describing what you lived and felt.
\end{grammarlens}

\subsection*{Your turn}
\begin{itemize}
\raggedright
  \item \emph{Say it:} \emph{fālij}
  \item \emph{Say it:} \emph{fālij}, once more
  \item \emph{Say it:} \emph{dādā ko fālij huā thā.} — My grandfather had had a stroke.
  \item \emph{Recall:} say the Urdu for chronic, then the Urdu for asthma, then say \emph{fālij} again
\end{itemize}

\begin{tcolorbox}[breakable,colback=teal!4,colframe=teal!35!black,title={Before you move on}]
What does \ur{فالج} mean? (\textquotedblleft{}Paralysis; a stroke\textquotedblright{}.) Say it once more.
\end{tcolorbox}
\section[\texorpdfstring{\ur{کینسر} (kainsar)}{کینسر (kainsar)}]{\ur{کینسر} (kainsar) — cancer}

\label{lesson:UR-C475-kainsar}



\begin{quote}
Before the new one: say the Urdu for asthma, then the Urdu for paralysis.
\end{quote}

\subsection*{You'll want to know: \ur{کینسر}}
\textbf{\ur{کینسر}} — \emph{kainsar} — \textquotedblleft{}cancer\textquotedblright{}.

\begin{quote}
\ur{کینسر} \ur{ایک} \ur{خطرناک} \ur{بیماری} \ur{ہے۔}
\end{quote}

\begin{quote}
\emph{kainsar ek khatarnāk bīmārī hai.}
\end{quote}

\begin{quote}
Cancer is a dangerous illness.
\end{quote}

Masculine.

\begin{cousinweb}
\textbf{\ur{کینسر}} (\emph{kainsar}) is borrowed from English \emph{cancer}.
\end{cousinweb}

\begin{grammarlens}[title={What you've built}]
One more word for describing what you lived and felt.
\end{grammarlens}

\subsection*{Your turn}
\begin{itemize}
\raggedright
  \item \emph{Say it:} \emph{kainsar}
  \item \emph{Say it:} \emph{kainsar}, once more
  \item \emph{Say it:} \emph{kainsar ek khatarnāk bīmārī hai.} — Cancer is a dangerous illness.
  \item \emph{Recall:} say the Urdu for asthma, then the Urdu for paralysis, then say \emph{kainsar} again
\end{itemize}

\begin{tcolorbox}[breakable,colback=teal!4,colframe=teal!35!black,title={Before you move on}]
What does \ur{کینسر} mean? (\textquotedblleft{}Cancer\textquotedblright{}.) Say it once more.
\end{tcolorbox}
\section[\texorpdfstring{\ur{مبتلا} (mubtalā)}{مبتلا (mubtalā)}]{\ur{مبتلا} (mubtalā) — afflicted (with)}

\label{lesson:UR-C475-mubtala}



\begin{quote}
Before the new one: say the Urdu for paralysis, then the Urdu for cancer.
\end{quote}

\subsection*{You'll want to know: \ur{مبتلا}}
\textbf{\ur{مبتلا}} — \emph{mubtalā} — \textquotedblleft{}afflicted (with)\textquotedblright{}.

\begin{quote}
\ur{وہ} \ur{دو} \ur{سال} \ur{سے} \ur{دمے} \ur{میں} \ur{مبتلا} \ur{ہے۔}
\end{quote}

\begin{quote}
\emph{voh do sāl se dame meṅ mubtalā hai.}
\end{quote}

\begin{quote}
He has suffered from asthma for two years.
\end{quote}

\begin{cousinweb}
\textbf{\ur{مبتلا}} (\emph{mubtalā}) is usually traced, through Persian, to Arabic \emph{mubtalā}, \textquotedblleft{}tried, afflicted\textquotedblright{}.
\end{cousinweb}

\begin{grammarlens}[title={What you've built}]
One more word for describing what you lived and felt.
\end{grammarlens}

\subsection*{Your turn}
\begin{itemize}
\raggedright
  \item \emph{Say it:} \emph{mubtalā}
  \item \emph{Say it:} \emph{mubtalā}, once more
  \item \emph{Say it:} \emph{voh do sāl se dame meṅ mubtalā hai.} — He has suffered from asthma for two years.
  \item \emph{Recall:} say the Urdu for paralysis, then the Urdu for cancer, then say \emph{mubtalā} again
\end{itemize}

\begin{tcolorbox}[breakable,colback=teal!4,colframe=teal!35!black,title={Before you move on}]
What does \ur{مبتلا} mean? (\textquotedblleft{}Afflicted (with)\textquotedblright{}.) Say it once more.
\end{tcolorbox}
